\section{The complete cover discrepancy generates the original zero-jet
ideal}\label{the-complete-cover-discrepancy-generates-the-original-zero-jet-ideal}

Independent mathematical derivation, 25 September 2026. Proof locators
\textbf{NCI0--NCI8}. This is a source-level calculation for the fixed
original receiver, with no parameter substitution, bounded zero search
or purity assumption. No remote publication is performed here.

\subsection{NCI0. Question, construction stage, and source
coverage}\label{nci0.-question-construction-stage-and-source-coverage}

The question is whether the complete family of prequotient cover
discrepancies already found in NHJ7 generates the entire original
summation ideal, in its original topology. This asks what the existing
defect actually records, rather than replacing the foundational
coefficient object by a different receiving space.

The support remains \(\tau\langle Z_1;\mathrm{no}\ Z_2\rangle\). The
scalar operations, logarithms, integers and analytic test spaces below
occur after the complete global arithmetic reconstruction. The separate
branch counters remain separate. Neither an arithmetic value, parity,
midpoint, metric nor the retracted addition is assigned to that support.
The current correction chain and the complete relevant user passages
were read for NHJ and retained in its source-use receipt; they remain
governing here.

Before deriving the ideal identity, the complete programme proof
\texttt{GLOBAL\_MELLIN\_SYNTHESIS.md}, S1--S7 and its later exact-image
notice, and the complete independent
\nolinkurl{GLOBAL_MELLIN_SYNTHESIS_REVIEW.md}, R1--R5, were read. The
global division estimate used below is their S3/R1--R3, with the
corrected fixed-height supremum. GSP0--GSP9 and GAP's exact source
comparison and original Gaussian estimate were already read in the
preceding independent derivation and were recalled. NHJ's complete
coefficient receiver and its exact prequotient cover formula are
retained. This is not a claim of new complete reading of all ancestor
literature.

Human provenance: S. Waleed Noor, \emph{A Hardy space analysis of the
Báez-Duarte criterion for the RH},
\href{https://arxiv.org/abs/1809.09577v4}{arXiv:1809.09577v4}, §§1--5,
original author TeX read in NHJ. The original source multiplier and
Hadamard comparison are those proved in S2 of
\texttt{GLOBAL\_MELLIN\_SYNTHESIS.md}, whose author-TeX routing and
precise historical source coverage remain in that proof. All new
ideal-generation and covariance conclusions are proved below and are not
attributed to Noor or Deligne.

\subsection{NCI1. The complete original coefficient system and its cover
discrepancy}\label{nci1.-the-complete-original-coefficient-system-and-its-cover-discrepancy}

Retain \[
\mathcal B=\left\{F\in\mathcal O(\mathbb C):
b_{A,N}(F)=\sup_{|\Re s|\le A}(1+|\Im s|)^N|F(s)|<\infty
\text{ for every integer }A,N\ge0\right\},
\tag{NCI1.1}
\] with precisely these Fréchet seminorms. The full ideal and quotient
are \[
\mathcal I=\{F\in\mathcal B:F^{(j)}(\rho)=0
\text{ for every actual nontrivial zero }\rho\text{ of original }\zeta,
\ 0\le j<m_\rho\},\qquad Q=\mathcal B/\mathcal I.
\tag{NCI1.2}
\] Each jet is continuous by Cauchy's formula on a small fixed disk.
Therefore \(\mathcal I\) is closed. No zero simplicity, zero separation
or restriction to finitely many zeros is part of this definition.

The actual source element is \[
F_0(s)=\frac{s(s-1)}8\pi^{-s/2}\Gamma(s/2)\zeta(s)\in\mathcal B,
\qquad F_0(0)=F_0(1)=\frac18.
\tag{NCI1.3}
\] Its full zero divisor is exactly NCI1.2's divisor. At its exceptional
points, \[
F_0(-1)=F_0(2)=\frac\pi{24},\qquad
F_0(-2k)=\frac{k(2k+1)(-1)^k\pi^k}{2\,k!}\zeta'(-2k)\ne0.
\tag{NCI1.4}
\] At every nontrivial zero its local germ remains \[
F_0(\rho+z)=z^{m_\rho}
\frac{(\rho+z)(\rho+z-1)}8\pi^{-(\rho+z)/2}
\Gamma((\rho+z)/2)\frac{\zeta(\rho+z)}{z^{m_\rho}}.
\tag{NCI1.5}
\] These full factors, and all their derivatives, are retained whenever
division or multiplicity occurs below. The auxiliary source element is
not a working replacement of the original zeta function.

Fix one \(t>0\) throughout, and define \[
g_t(s)=e^{ts^2},\qquad U_nF(s)=n^{1-s}F(s),
\] \[
\boxed{\delta_{n,t}(s)=8e^{ts^2}F_0(s)
\frac{n-n^{1-s}}s,\qquad n=2,3,\ldots.}
\tag{NCI1.6}
\] The entire quotient at zero is \(n\log n\). On each fixed strip, the
quotient is bounded outside the unit disk and bounded inside by its
entire removable extension. Gaussian multiplication makes it rapidly
decreasing on that strip. Consequently \[
\delta_{n,t}\in F_0\mathcal B\subset\mathcal I.
\tag{NCI1.7}
\] The endpoint values are not zero conditions: \[
\delta_{n,t}(0)=n\log n,\qquad
\delta_{n,t}(1)=e^t(n-1),
\] \[
\delta_{n,t}(-2k)=8e^{4tk^2}F_0(-2k)
\frac{n-n^{1+2k}}{-2k}.
\tag{NCI1.8}
\] The full expression at each such point is retained. Only the actual
nontrivial-zero jets are imposed by membership in \(\mathcal I\).

\subsection{NCI2. Global division and the closure of the original
principal
image}\label{nci2.-global-division-and-the-closure-of-the-original-principal-image}

The already proved global source estimates, with multiplicities
retained, are \[
F_0(s)=e^{a+bs}\prod_\rho\left(1-\frac{s}{\rho}\right)e^{s/\rho},
\quad e^a=\frac18,\quad b=\frac{F_0'(0)}{F_0(0)},
\] \[
\sum_{|\rho|\le R}m_\rho\le C(R+2)^{3/2}.
\tag{NCI2.1}
\] The product in the first line repeats each zero according to its
order; the second line counts distinct locations with those orders. S2
derives these from the full original Mellin kernel and the original
genus-one factorization, without changing its exponential factor.

Here is the division estimate, including clusters. For
\(F\in\mathcal I\), Taylor division at every actual zero extends
\(H=F/F_0\) to an entire function. On a large scale \(R\), take the
union of open disks of radius \(R^{-2}\) around the zeros of modulus at
most \(8R\). Their total radii are at most \(CR^{-1/2}\), by NCI2.1.
Each connected component has diameter at most twice the sum of its disk
radii, hence less than one for sufficiently large \(R\).

For \(R/2\le|s|\le3R\) outside these disks, the finite Hadamard factors
have \[
|1-s/\rho|\ge(8R^3)^{-1}\quad(|\rho|\le8R).
\tag{NCI2.2}
\] Their logarithms sum to at least \(-CR^{3/2}\log R\). The retained
exponentials contribute at least \[
-3R\sum_{|\rho|\le8R}m_\rho/|\rho|\ge-CR^{3/2}
\]. For the tail, \(|s/\rho|\le3/8\), and
\(\log|(1-w)e^w|\ge-\frac85|w|^2\). Partial summation of NCI2.1 gives
\(\sum_{|\rho|>8R}m_\rho/|\rho|^2\le CR^{-1/2}\), so this tail
contributes at least \(-CR^{3/2}\). The leading exponential contributes
at least \(-|a|-3|b|R\). Therefore the full product obeys \[
|F_0(s)|^{-1}\le\exp(CR^{3/2}\log R)
\tag{NCI2.3}
\] on the stated complement. Every factor and multiplicity has been
counted.

For \(|\Re s|\le A\), \(R\le|s|\le2R\), outside the disks, multiply
NCI2.3 by the strip bound for \(F\). Inside a disk component, its
boundary lies in \(|\Re s|\le A+1\) and \(R/2\le|s|\le3R\), and outside
the open disk union. Apply that boundary estimate to the entire function
\(H\), and then the maximum-modulus principle on the compact component
closure. Smoothness or simple connectivity of its boundary is
unnecessary. This also handles multiple or arbitrarily close zeros,
since all local singularities of \(H\) have already been removed by the
full multiplicity hypothesis. Taking dyadic \(R\) and absorbing the
remaining compact region proves \[
\sup_{|x|\le A}|H(x+iy)|
\le C_{A,F}\exp\{C(|y|+2)^{3/2}\log(|y|+2)\}.
\tag{NCI2.4}
\] The supremum is at a fixed imaginary coordinate, exactly as corrected
in the independent source review.

For every \(\varepsilon>0\), the entire quotient \[
H_\varepsilon(s)=e^{\varepsilon s^2}\frac{F(s)}{F_0(s)}
\tag{NCI2.5}
\] belongs to \(\mathcal B\): on the strip its modulus is bounded by
\(e^{\varepsilon A^2-\varepsilon y^2}\) times NCI2.4, and the negative
quadratic dominates its subquadratic exponent and every polynomial
weight. Its exact image under multiplication is \[
F_0H_\varepsilon=e^{\varepsilon s^2}F.
\tag{NCI2.6}
\] This is the unshifted Gaussian used in NHJ and GAP. The earlier S4
approximation used \(e^{\varepsilon(s-1/2)^2}\); the same proved
division bound validates NCI2.5 directly, and no source formula is
silently identified with the centered one.

For \(0<\varepsilon\le1\), integrating \(s^2e^{v s^2}\) from zero to
\(\varepsilon\) gives the exact source estimate \[
b_{A,N}((e^{\varepsilon s^2}-1)F)
\le\varepsilon e^{A^2}(A^2+1)b_{A,N+2}(F).
\tag{NCI2.7}
\] Hence NCI2.6 converges to \(F\) in the original topology. Combining
this with NCI1.7 and closedness of \(\mathcal I\) proves \[
\boxed{\overline{F_0\mathcal B}^{\mathcal B}=\mathcal I.}
\tag{NCI2.8}
\] The algebraic equality is false: \(F_0\in\mathcal I\), but
\(F_0=F_0G\), \(G\in\mathcal B\), would force \(G=1\) on the nonzero
open set of \(F_0\), hence everywhere. The constant one is not in
\(\mathcal B\). Thus \[
F_0\mathcal B\subsetneq\mathcal I,
\tag{NCI2.9}
\] with an explicit missing element, while the closure is exactly
NCI2.8. No principal-ideal equality has been assumed.

\subsection{NCI3. The entire parameter family remembers all recovered
integers}\label{nci3.-the-entire-parameter-family-remembers-all-recovered-integers}

Let \(\Lambda\in\mathcal B'\) annihilate \(\delta_{n,t}\) for every
\(n\ge2\). Define \[
K_{t,w}(s)=8g_t(s)F_0(s)
\frac{e^w-e^{(1-s)w}}s,
\qquad L(w)=\Lambda(K_{t,w}).
\tag{NCI3.1}
\] The numerator vanishes at \(s=0\); its quotient there equals
\(we^w\). The exact integral expression \[
K_{t,w}(s)=8we^w g_t(s)F_0(s)\int_0^1e^{-\theta sw}\,d\theta
\tag{NCI3.1a}
\] follows by integrating the exponential in \(\theta\). It retains the
full original expression and verifies directly that \(K_{t,w}(0)=we^w\).
Thus this is an entire source function for every \(w\), not a
meromorphic test.

It is also an entire \(\mathcal B\)-valued family. For a precise bound,
write \(s=x+iy\), \(w=a+ib\). Outside \(|s|<1\) on a fixed real strip,
the first numerator term is bounded after Gaussian multiplication by
\(e^{tA^2+|a|-ty^2}\). The second is bounded by \[
e^{tA^2+(1+A)|a|}(1+|y|)^N e^{-ty^2+|b||y|}
\tag{NCI3.2}
\] in its weighted supremum, multiplied by the fixed bounded-strip
factor \(8F_0\). The inequality \(-ty^2+|b||y|\le-(t/2)y^2+b^2/(2t)\)
absorbs every polynomial factor. On the inner unit disk the maximum
principle for the entire divided function on \(|s|=2\) bounds it by a
constant times \(e^{3|w|}\). Therefore \[
b_{A,N}(K_{t,w})\le C_{A,N,t}\exp(C_{A,t}|w|+|w|^2/(2t)).
\tag{NCI3.3}
\] Parameter derivatives add powers of \(1-s\) and the same type of
bound on compact \(w\)-sets; Taylor remainders are controlled as well.
This proves source-valued holomorphy and consequently entire scalar
holomorphy of \(L\), with \(|L(w)|\le C\exp(C(1+|w|^2))\).

Exactly, with no finite-range approximation, \[
K_{t,\log n}=\delta_{n,t}\quad(n\ge2),\qquad K_{t,0}=0.
\tag{NCI3.4}
\] Thus \(L\) vanishes at every \(\log n\), \(n\ge1\). If it were
nonzero, remove its finite-order zero at zero and apply Jensen's
formula. On a disk of radius \(2R\), its quadratic exponential growth
permits at most \(O(R^2)\) zeros in radius \(R\), because each such zero
contributes at least \(\log2\) to Jensen's sum. But the distinct zeros
\(\log2,\ldots,\log\lfloor e^R\rfloor\) give exponentially many. This
contradiction proves \(L\equiv0\).

Differentiate the full source expression before taking a quotient: \[
(\partial_w-1)K_{t,w}(s)
=8g_t(s)F_0(s)e^{(1-s)w}.
\tag{NCI3.5}
\] The sign is positive: subtracting the undifferentiated second
exponential leaves \(s e^{(1-s)w}\) in the numerator. Applying
\(\Lambda\), dividing the nonzero scalar \(8e^w\), and putting \(v=-w\)
yields \[
\boxed{\Lambda(F_0g_t e^{vs})=0\quad(v\in\mathbb C).}
\tag{NCI3.6}
\]

\subsection{NCI4. Exact original Mellin synthesis of the
annihilator}\label{nci4.-exact-original-mellin-synthesis-of-the-annihilator}

Multiplication by the fixed \(F_0\) is continuous on \(\mathcal B\),
since \(b_{A,N}(F_0F)\le b_{A,0}(F_0)b_{A,N}(F)\). Therefore
\(\ell(F)=\Lambda(F_0F)\) is a continuous functional on the original
source. Equation NCI3.6 says it annihilates every translated Gaussian
\(g_t e^{vs}\). We now prove that it annihilates all of \(\mathcal B\),
with the exact source constants.

The original source is \[
A=\{a\in C^\infty(\mathbb R_{>0}):
\sup_{u>0}(u^N+u^{-N})|(u\partial_u)^j a(u)|<\infty
\text{ for all }N,j\ge0\},
\] \[
\Theta a(s)=\frac12\int_0^\infty a(u)u^s\frac{du}{u},\qquad
\Theta^{-1}H(u)=\frac{u^{-1/2}}\pi
\int_{\mathbb R}H(1/2+iy)u^{-iy}\,dy.
\tag{NCI4.1}
\] These exact inverse topological maps were proved in the source and
retained in GAP3/NHJ1. For arbitrary \(F\in\mathcal B\) and real
\(u>t\), let \(H=e^{(u-t)s^2}F\in\mathcal B\), \(a_H=\Theta^{-1}H\). Its
full half-Mellin integral gives \[
e^{us^2}F(s)=\frac12\int_{\mathbb R}
a_H(e^v)g_t(s)e^{sv}\,dv.
\tag{NCI4.2}
\] It converges in every \(\mathcal B\) seminorm: the source kernel has
\(b_{A,N}(g_t e^{sv})\le C_{A,N,t}e^{A|v|}\) for real \(v\), while
\(|a_H(e^v)|\le C_Me^{-M|v|}\) for every \(M\). Choose \(M>A+1\).
Continuity permits passing \(\ell\) through this locally convex
integral, and NCI3.6 gives \(\ell(e^{us^2}F)=0\) for every real \(u>t\).

For fixed \(F\), this scalar expression is holomorphic for \(\Re u>0\).
Indeed, writing \(u=p+iq\) and \(s=x+iy\), its Gaussian logarithmic
modulus is \(p(x^2-y^2)-2qxy\). On compact subsets of that parameter
half-plane the negative quadratic uniformly absorbs the linear term and
every factor \(s^{2j}\) from parameter derivatives. This proves
source-valued holomorphy by dominated Taylor remainders in every source
seminorm. The identity theorem extends the zero to the whole half-plane.
Finally NCI2.7 and continuity let real \(u\downarrow0\), giving
\(\ell(F)=0\). Thus \[
\Lambda(\delta_{n,t})=0\ (\forall n\ge2)
\quad\Longrightarrow\quad
\Lambda(F_0\mathcal B)=0.
\tag{NCI4.3}
\] The converse holds from NCI1.7.

\subsection{NCI5. The full closed ideal is generated by the cover
discrepancies}\label{nci5.-the-full-closed-ideal-is-generated-by-the-cover-discrepancies}

Put \(S_t=\operatorname{span}_{\mathbb C}\{\delta_{n,t}:n\ge2\}\), where
the span means finite linear combinations before closure. By NCI1.7,
\(S_t\subset F_0\mathcal B\subset\mathcal I\). Equations NCI4.3 and
NCI2.8 identify their continuous annihilators: \[
S_t^\perp=(F_0\mathcal B)^\perp=\mathcal I^\perp
\quad\text{in }\mathcal B'.
\tag{NCI5.1}
\] If \(\overline{S_t}\) were a proper closed subspace of
\(\mathcal I\), choose an element of
\(\mathcal I\setminus\overline{S_t}\). Hahn--Banach separation in the
Hausdorff locally convex space \(\mathcal B\) gives a continuous
functional vanishing on \(\overline{S_t}\) but not on that element.
Equation NCI5.1 contradicts this. Therefore \[
\boxed{\overline{\operatorname{span}\{\delta_{n,t}:n\ge2\}}^{\mathcal B}
=\overline{F_0\mathcal B}^{\mathcal B}
=\mathcal I\qquad\text{for the fixed }t>0.}
\tag{NCI5.2}
\] This is the full original zero-jet ideal, not a radical ideal
omitting multiplicities, a value-only ideal, or a finite-height
assertion. The proof needs the complete family of recovered integers. A
single covariance equation has only the single annihilator
\(\ker(\Lambda\mapsto\Lambda(\delta_{n,t}))\); it has not been
substituted for NCI5.1.

In particular the induced map \[
\mathcal B/\overline{S_t}\longrightarrow Q,
\qquad [F]_{\overline{S_t}}\longmapsto[F]_{\mathcal I}
\tag{NCI5.3}
\] is a topological isomorphism with inverse given by the same
representative. Well-definedness, injectivity and surjectivity follow
from the proved equality of the two kernels; the quotient topologies
coincide because their original source map and kernel coincide. Thus the
closed covariance discrepancy constructs the already specified source
quotient exactly.

The original arithmetic-source version retains the full inverse
transform: \[
c_{n,t}(u)=\frac{u^{-1/2}}\pi
\int_{\mathbb R}\delta_{n,t}(1/2+iy)u^{-iy}\,dy.
\tag{NCI5.4}
\] Each integral is an element of \(A\), and the proved inverse
topological isomorphism gives \[
\overline{\operatorname{span}\{c_{n,t}:n\ge2\}}^A
=\Theta^{-1}\mathcal I.
\tag{NCI5.5}
\] The already proved exact-image theorem identifies this last space
with the original \(J=\Sigma S\), \(\Sigma f=2\sum_{k\ge1}f(ku)\),
including its even-Schwartz endpoint constraints. This invokes that
proved source identification rather than claiming a new Schwartz
inversion theorem here. Thus the actual original ideal appears in both
source presentations with their exact map and constants.

\subsection{NCI6. All-cover covariance characterizes precisely the
original quotient
dual}\label{nci6.-all-cover-covariance-characterizes-precisely-the-original-quotient-dual}

Retain the original corrected coefficient tests of NHJ: \[
\psi_0=\frac{-1+8F_0}s,\qquad
\psi_m=\frac{m^{1-s}-(m+1)^{1-s}+8F_0}s\quad(m\ge1),
\] \[
\mathcal H_t^{\mathcal B}\Lambda(z)
=\sum_{m\ge0}\overline{\Lambda(g_t\psi_m)}z^m.
\tag{NCI6.1}
\] All divisions are the entire removable ones already proved. Its
continuous injection \(\mathcal B'_\beta\to\mathcal O(\mathbb D)\) is
NHJ2--NHJ4. The raw analytic cover adjoint sums consecutive coefficient
blocks: \((\mathscr W_n^*f)_m=\sum_{a=0}^{n-1}f_{nm+a}\). Finite
telescoping gives the full prequotient identity \[
\mathscr W_n^*\mathcal H_t^{\mathcal B}\Lambda
-\mathcal H_t^{\mathcal B}U_n'\Lambda
=\frac{\overline{\Lambda(\delta_{n,t})}}{1-z}.
\tag{NCI6.2}
\] No zero or boundary term is omitted. Because \(1/(1-z)\) is a nonzero
holomorphic function on the disk, NCI5.1 proves \[
\boxed{
\mathscr W_n^*\mathcal H_t^{\mathcal B}\Lambda
=\mathcal H_t^{\mathcal B}U_n'\Lambda\quad\forall n\ge2
\iff\Lambda\in\mathcal I^\perp
\iff\Lambda=\lambda\circ q\text{ for a unique }\lambda\in Q'.}
\tag{NCI6.3}
\] The last equivalence is exact for the original quotient topology: a
continuous functional vanishing on the kernel descends continuously by
its definition. Thus the original full-jet quotient dual is
\emph{derived} as the maximal source subspace on which this fixed
receiver intertwines every raw arithmetic cover. It has not been imposed
as an extra definition of covariance.

This can also be written as an exact detector of the remaining ideal
directions. Restrict a functional to \(\mathcal I\), and then evaluate
on the discrepancy family: \[
\mathcal B'\xrightarrow{\mathrm{res}}\mathcal I'
\xrightarrow{\mathfrak d_t}\mathbb C^{\{2,3,\ldots\}},
\qquad\mathfrak d_t(\mu)=(\mu(\delta_{n,t}))_{n\ge2}.
\tag{NCI6.4}
\] Restriction is surjective as a map of vector spaces by continuous
Hahn--Banach extension from the closed subspace \(\mathcal I\); it is
strongly continuous because bounded subsets of that subspace are bounded
in \(\mathcal B\). Its kernel is precisely \(q'Q'\). The second map is
injective by NCI5.2 and continuous into the product topology since each
coordinate is evaluation on a singleton bounded set. No surjectivity
onto arbitrary sequences, openness of restriction in strong topologies,
or inverse continuity of this detector is assumed.

\subsection{NCI7. Exact all-cover stability of the Hardy
domain}\label{nci7.-exact-all-cover-stability-of-the-hardy-domain}

The actual prequotient Hilbert domain, with no hidden boundary
assumption, is \[
\mathfrak D_t^{\mathcal B}
=\left\{\Lambda\in\mathcal B'_\beta:
\sum_{m\ge0}|\Lambda(g_t\psi_m)|^2<\infty\right\}.
\tag{NCI7.1}
\] Its closed injective graph and complete graph topology were proved in
NHJ6. For \(\Lambda\) in this domain,
\(\mathscr W_n^*\mathcal H_t^{\mathcal B}\Lambda\) is in \(H^2\), since
the restriction is Noor's bounded raw \(W_n^*\). The right side of
NCI6.2 belongs to \(H^2\) precisely when its scalar coefficient is zero,
because all coefficients of \(1/(1-z)\) equal one. Hence \[
U_n'\Lambda\in\mathfrak D_t^{\mathcal B}
\iff\Lambda(\delta_{n,t})=0
\quad(\Lambda\in\mathfrak D_t^{\mathcal B}).
\tag{NCI7.2}
\] Using all recovered integers and NCI5.1 now proves \[
\boxed{
\{\Lambda\in\mathfrak D_t^{\mathcal B}:
U_n'\Lambda\in\mathfrak D_t^{\mathcal B}\ \forall n\ge2\}
=\mathfrak D_t^{\mathcal B}\cap\mathcal I^\perp
=q'\mathfrak D_t^Q.}
\tag{NCI7.3}
\] This answers the all-cover source-domain question exactly. It is
stronger than a family of separate sufficient conditions: every vector
with the complete raw-cover stability property must annihilate
\emph{all} original summation-source modes and every full multiplicity
constraint encoded by \(\mathcal I\).

The right side is invariant under every \(U_n'\). To prove this
directly, \(U_n\mathcal I\subset\mathcal I\) because its entire
multiplier preserves the full zero orders, and NCI7.2 preserves the
Hilbert domain; apply both properties. Thus it is the largest linear
subspace of \(\mathfrak D_t^{\mathcal B}\) stable under all the raw
cover transposes. If another such subspace is given, each of its vectors
satisfies the left side of NCI7.3 and therefore lies in this one. No
invariance under the inverse source action is inferred.

On this actual quotient domain the degree identities remain \[
W_n^*W_n=nI,\qquad W_nW_n^*=nP_n,
\tag{NCI7.4}
\] where \(P_n\) averages consecutive coefficient blocks. For the
transported positive form of NHJ6, the exact remaining defect is \[
nB_t(\lambda,\mu)-B_t(U_n'\lambda,U_n'\mu)
=n\langle(I-P_n)\mathcal H_t\mu,(I-P_n)\mathcal H_t\lambda\rangle.
\tag{NCI7.5}
\] The generation of the source ideal removes the scalar pole in NCI6.2
through its specified quotient. It does not remove this separate
projection defect, prove a global inverse of \(W_n^*\), or identify the
transported positive form with the original Weil pairing.

\subsection{NCI8. The source-level conclusion and its actual
scope}\label{nci8.-the-source-level-conclusion-and-its-actual-scope}

The fixed original cover discrepancy generates the complete original
zero-jet ideal after closure. Both stages have been proved: global
Gaussian division gives \(\overline{F_0\mathcal B}=\mathcal I\), while
the full integer-logarithm sampling and original half-Mellin integral
give
\(\overline{\operatorname{span}\delta_{n,t}}=\overline{F_0\mathcal B}\).
The source principal image is strictly smaller algebraically, and its
topology is essential.

The next calculation was not left as a possible application. NCI6--NCI7
apply the identity to the actual prequotient receiver, proving that its
full cover equivariance, and its full forward-cover Hardy-domain
stability, select exactly the original quotient dual. In particular the
quotient is recovered from the measured complete defect rather than
substituted merely because it is convenient.

The actual specialization \(\mathcal R=Q/N_O\) still has its further
source kernel \(N_O\); NCI6.3 selects \(Q'\), not just \(N_O^\perp\).
All original nontrivial zero jets are retained at this stage, including
off-line and higher jets. Therefore this ideal identity does not prove
\(\mathcal R=0\), remove the actual specialization coefficient
\(d_r(\rho)\), or supply Deligne's weight-separated vanishing. The
additional operation needed is still a proved geometric return on the
actual source with its exact degree and adjoint. NCI7.5 supplies the
precise remaining receiving defect that such a calculation must actually
evaluate, without assuming its vanishing.

One further consequence strengthens the generating-set statement without
changing the original receiver or its fixed parameter. Let
\(\mathscr S\subset\{2,3,\ldots\}\) be a set of recovered integer cover
degrees satisfying \[
\limsup_{R\to\infty}
\frac{\#\{n\in\mathscr S:\log n\le R\}}{R^2}=\infty.
\tag{NCI8.1}
\] Then \[
\overline{\operatorname{span}\{\delta_{n,t}:n\in\mathscr S\}}^{\mathcal B}
=\mathcal I.
\tag{NCI8.2}
\] To prove this, repeat NCI3 with a functional annihilating that
specified family. Its scalar entire function still has the same
quadratic growth bound. If nonzero, Jensen bounds its zeros in radius
\(R\) by \(CR^2\), contradicting NCI8.1. Every later step NCI4--NCI5 is
unchanged. Consequently NCI6.3 and NCI7.3 hold with the family
\(n\in\mathscr S\) on their left sides. For a fixed positive integer
\(d\), the complete family \(\mathscr S=\{k^d:k\ge2\}\) satisfies the
hypothesis because its counting numerator is
\(\lfloor e^{R/d}\rfloor-1\). In particular, all square-degree covers
already determine this same original ideal and the same all-cover-stable
Hardy source domain. This uses an infinite generating family within the
completely reconstructed arithmetic; it supplies no construction of the
integers from a finite selected spectrum.
